\documentclass[../main.tex]{subfiles}
\begin{document}

\section{Реализация системы}

\subsection{Структура проекта}

Исходный код сервера организован по слоям, приведённым
в таблице~\ref{tab:backend-layers}.

\begin{table}[H]
  \centering
  \small
  \begin{tabular}{@{}>{\raggedright\arraybackslash}p{2.4cm}>{\raggedright\arraybackslash}p{6.7cm}>{\raggedright\arraybackslash}p{5.0cm}@{}}
    \toprule
    Слой & Классы & Ответственность \\
    \midrule
    \field{controller} & \field{DbController}, \field{RowsController},
      \field{JoinController} & REST API \\
    \field{service} & \field{SchemaLoader}, \field{DiffService},
      \field{RowsService}, \field{FederationService} и др. & Сбор схемы,
      сравнение, выборка строк, федеративные запросы \\
    \field{repository} & \field{InstanceRepository},
      \field{CatalogRepository}, \field{QueryHistoryRepository} &
      Доступ к таблицам \field{sc\_*} \\
    \field{model} & \field{Catalog}, \field{Schema}, \field{Table},
      \field{Column}, \field{Constraint}, \field{Index} & Объектная
      модель схемы и результатов \\
    \field{util} & \field{SqlQueries}, \field{CryptoUtil},
      \field{RowsUtil} & SQL-константы, шифрование \\
    \bottomrule
  \end{tabular}
  \caption{Слои серверной части}
  \label{tab:backend-layers}
\end{table}

Клиентская часть построена из компонентов: \field{InstanceDock}
с формой подключения и деревом схем, \field{SchemaTree},
\field{RowsExplorer} и \field{RowDiffView} для работы с данными,
\field{SchemaDiffDrawer} и \field{DiffViewer} для сравнения схем,
\field{JoinBuilder} и \field{QueryResultView} для федеративных
запросов. Состояние вынесено в runes-модули
\field{instances.svelte.ts} и \field{join.svelte.ts}; модуль
\field{api.ts} инкапсулирует обращения к REST.

\subsection{Сбор и сохранение метаданных}

Загрузка схемы начинается POST-запросом к ресурсу инстанса
(\field{/load}). Обработчик проверяет существование инстанса,
расшифровывает пароль и выбирает коннектор
по типу СУБД. Коннектор \field{PostgresConnector} читает
\field{information\_schema} и \field{pg\_catalog}, а
\field{MariaDbConnector}~--- \field{information\_schema} MariaDB.
Оба коннектора работают через соединение <<только чтение>>, поэтому
сбор метаданных не блокирует таблицы источника и не требует его
остановки.

Собранный каталог сохраняется в рамках транзакции в два этапа.
Сначала записываются схемы, таблицы, колонки и индексы. Затем
сохраняются ограничения: для внешнего ключа выполняется поиск
целевой таблицы в уже сохранённом каталоге и заполняется
\field{ref\_table\_id}. Если целевая таблица не найдена, внешний
ключ пропускается с записью в журнал, чтобы не нарушать целостность
каталога. После успешного сохранения каталог становится доступен
для последующих запросов.

\subsection{REST API}

Все эндпоинты расположены под префиксом \path{/api}. Основные
операции приведены в таблице~\ref{tab:api}. Запросы возвращают
полный JSON-ответ и завершаются сразу после отправки данных.

{\footnotesize
\begin{longtable}{@{}p{1.25cm}p{6.0cm}p{7.5cm}@{}}
  \caption{Основные эндпоинты REST API}
  \label{tab:api}\\
  \toprule
  Метод & Путь & Назначение \\
  \midrule
  \endfirsthead
  \toprule
  Метод & Путь & Назначение \\
  \midrule
  \endhead
  \bottomrule
  \endfoot
  GET & \path{/api/instances} & Список подключённых источников \\
  GET & \path{/api/instances/{id}} & Источник по идентификатору \\
  POST & \path{/api/instances} & Создание источника \\
  PATCH & \path{/api/instances/{id}/password} & Смена пароля \\
  PATCH & \path{/api/instances/{id}/alias} & Смена alias для федерации \\
  POST & \path{/api/instances/aliases/backfill} & Проставить alias
    существующим источникам \\
  POST & \path{/api/instances/{id}/load} & Собрать и сохранить схему \\
  POST & \path{/api/instances/{id}/check} & Сравнить сохранённую схему
    с текущей \\
  GET & \path{/api/instances/{id}/schema} & Каталог источника \\
  POST & \path{/api/instances/compare} & Сравнить две схемы \\
  GET & \path{/api/instances/{id}/tables/{schema}/{table}/rows} &
    Страница строк таблицы \\
  POST & \path{/api/join} & Выполнить федеративный запрос \\
  GET & \path{/api/join/history} & История федеративных запросов \\
\end{longtable}
}

\subsection{Федеративные запросы}

Федеративный запрос формируется клиентом в виде структуры
\field{\{tables, select, maxRows, offset, timeoutMs\}} и отправляется
в \path{POST /api/join}. Сервер компилирует структуру в SQL, разбирает
его через JSqlParser, определяет источники и нужные колонки, а затем
выполняет чтение подтаблиц по JDBC. Прочитанные строки загружаются
во временные таблицы хранилища PostgreSQL, где и выполняется
соединение. Такой подход позволяет соединять данные разных СУБД
без переноса всех данных в одно постоянное хранилище:
материализуются только строки и колонки, необходимые для конкретного
запроса.

Для ограничения нагрузки на источники и хранилище установлены
лимиты: до 1000 строк в ответе по умолчанию и до 100~000
максимально, время выполнения до 30~с по умолчанию и до 300~с
максимально, до 500~000 строк из одного источника. Результат
истории сохраняется в таблице \field{sc\_query\_history}.

\end{document}
